\documentclass[10pt,a4paper,ssfamily]{exam}
\usepackage[brazil]{babel}
\usepackage[numbers,sort&compress]{natbib}
\usepackage[utf8]{inputenc}
\renewcommand{\baselinestretch}{1.0}
\usepackage{epsfig}
\usepackage{mhchem}
\usepackage{graphicx}
\usepackage{makeidx}
\usepackage{multicol}
\usepackage{multirow}
\usepackage{amssymb}
\usepackage{amsmath}
\DeclareMathOperator{\sen}{sen}
\newcommand{\listfont}{\bf\fontencoding{OT1}\sffamily\large}
\newcommand{\titlesfont}{\fontencoding{OT1}\sffamily\large}
\renewcommand{\baselinestretch}{1.0}
\newcommand{\Mg}{Mg$^{2+}$}
\renewcommand{\t}{\textrm}
\newcommand{\cc}[1]{[\textrm{#1}]}
\usepackage{enumitem}
\newcommand{\bmat}{\left[\begin{array}{cc}}
\newcommand{\emat}{\end{array}\right]}
\newcommand{\nomera}{
  \noindent
  \begin{tabular}{|p{0.7\textwidth}|p{0.3\textwidth}|}
  Nome: & RA: \\
  \hline
  \end{tabular}
}

\newcommand{\qini}{\begin{enumerate}[resume]\item}
\newcommand{\qend}{\end{enumerate}}

\newcommand{\oC}{$^\circ$C}
\newcommand{\kcalmol}{~kcal~mol$^{-1}$}
\newcommand{\moll}{~mol~L$^{-1}$}
\newcommand{\molkg}{~mol~kg$^{-1}$}
\newcommand{\gl}{~g~L$^{-1}$}
\newcommand{\e}[1]{$\times 10^{#1}$}

\newcommand{\eps}{\varepsilon}

\newcommand{\Resposta}[1]{\vspace{0.5cm}\noindent{\bf Resposta:}\\ #1
\begin{center}\rule{\textwidth}{0.4pt}\end{center}}
\renewcommand{\Resposta}[1]{}
\newcommand{\resposta}[1]{\vspace{0.5cm}\noindent{\bf Resposta:}\\ #1
\begin{center}\rule{\textwidth}{0.4pt}\end{center}}


\begin{document}
\sffamily
\pagestyle{empty}
\nomera

\begin{center}

{\bf QG101 - Química Geral}

{\bf Aula 2 - Estequiometria}\\
\underline{Justifique suas respostas.}
\end{center}

\begin{center}{\bf Questões}\end{center}

% ---------------------------------------------------------------
\qini
A \textbf{termita} é uma reação altamente exotérmica usada em soldagem
e em incêndios industriais. A equação não balanceada é:

\[
  \ce{Al(s) + Fe2O3(s) -> Al2O3(s) + Fe(l)}
\]

\begin{enumerate}
  \item Balanceie a equação química acima.
  \item Quantos gramas de \ce{Al} ($M = 26{,}98~\text{g/mol}$) são
        necessários para reagir completamente com $160{,}0~\text{g}$
        de \ce{Fe2O3} ($M = 159{,}7~\text{g/mol}$)?
  \item Qual a massa de ferro líquido (\ce{Fe},
        $M = 55{,}85~\text{g/mol}$) produzida nessa reação?
\end{enumerate}

\vspace{2cm}
\qend

% ---------------------------------------------------------------
\qini
O \textbf{processo Haber-Bosch} produz amônia para fertilizantes:
\[
  \ce{N2(g) + 3 H2(g) -> 2 NH3(g)}
\]
Em um reator industrial, misturam-se $56{,}0~\text{g}$ de \ce{N2}
($M = 28{,}02~\text{g/mol}$) com $9{,}00~\text{g}$ de \ce{H2}
($M = 2{,}016~\text{g/mol}$).
\begin{enumerate}
  \item Determine o número de mols de cada reagente.
  \item Identifique o \textbf{reagente limitante}, mostrando o cálculo.
  \item Calcule a massa teórica de \ce{NH3} ($M = 17{,}03~\text{g/mol}$)
        produzida e a massa do reagente em excesso que sobra.
\end{enumerate}

\vspace{2cm}
\qend

% ---------------------------------------------------------------
\qini
Em um laboratório de análise de águas, realiza-se uma \textbf{titulação}
para determinar a concentração de ácido clorídrico (\ce{HCl}) em uma
amostra de efluente industrial. A reação é:
\[
  \ce{HCl(aq) + NaOH(aq) -> NaCl(aq) + H2O(l)}
\]
Para neutralizar $50{,}0~\text{mL}$ da amostra foram necessários
$32{,}4~\text{mL}$ de \ce{NaOH} $0{,}150~\text{mol/L}$.
\begin{enumerate}
  \item Calcule o número de mols de \ce{NaOH} utilizados e, pela
        estequiometria, o número de mols de \ce{HCl} na amostra.
  \item Calcule a concentração molar de \ce{HCl} na amostra.
  \item Converta essa concentração para g/L
        ($M_{\ce{HCl}} = 36{,}46~\text{g/mol}$).
        O resultado é extensivo ou intensivo? Justifique.
\end{enumerate}

\vspace{2cm}
\qend

\end{document}
